\chapter{L'apparato digerente: cavità orale, faringe ed esofago}

% ---------------------------------------------------------------------------
\section{Il peritoneo e l'organizzazione della cavità addominale}

La cavità addominale è rivestita da una membrana detta \textbf{peritoneo},
costituita da due strati: uno \textbf{viscerale}, che avvolge gli organi, e
uno \textbf{parietale}, che riveste la superficie interna della parete
corporea; i due strati sono in continuità tra loro. All'interno della
cavità addominale sono racchiusi numerosi organi, che in base al loro
rapporto con il peritoneo si distinguono in tre gruppi: gli organi
\textbf{intraperitoneali} (stomaco, fegato, ileo), completamente avvolti
dal peritoneo viscerale; gli organi \textbf{retroperitoneali} (reni,
ureteri, aorta addominale), posti dietro il peritoneo parietale; e gli
organi \textbf{retroperitoneali secondari} (pancreas e parte del duodeno),
che durante lo sviluppo perdono il loro rivestimento peritoneale originario
e si addossano alla parete posteriore dell'addome.

% ---------------------------------------------------------------------------
\section{Generalità e funzioni dell'apparato digerente}

L'apparato digerente è costituito da un insieme di organi che processa il
cibo, estrae i nutrienti e ne elimina i residui. Il canale digerente e gli
organi accessori ad esso annessi cooperano per svolgere sei funzioni
principali:

\begin{enumerate}
  \item \textbf{Ingestione}: assunzione selettiva di cibo.
  \item \textbf{Digestione}: frammentazione meccanica e chimica del cibo in
    forma utilizzabile dall'organismo.
  \item \textbf{Secrezione}: acidi, enzimi e tamponi vengono secreti dalla
    parete del canale digerente o dagli organi accessori annessi.
  \item \textbf{Assorbimento}: passaggio delle molecole nutrienti nel
    sangue e nella linfa.
  \item \textbf{Compattazione}: progressiva disidratazione del materiale
    non digeribile e dei rifiuti organici.
  \item \textbf{Defecazione}: eliminazione del materiale fecale.
\end{enumerate}

Il rivestimento del tratto digerente svolge inoltre un ruolo difensivo
contro l'effetto corrosivo di acidi ed enzimi digestivi, contro gli stress
meccanici e contro gli agenti patogeni residenti nel canale digerente o
ingeriti con il cibo.

I principali organi che compongono l'apparato digerente, con la relativa
funzione primaria, sono: la cavità orale con denti e lingua (trattamento
meccanico e umidificazione del cibo, mescolamento con le secrezioni
salivari); le ghiandole salivari (secrezione di fluido lubrificante
contenente enzimi che scindono i carboidrati); la faringe (i muscoli
faringei spingono il materiale nell'esofago); l'esofago (trasporto del
materiale allo stomaco); lo stomaco (scissione chimica dei materiali
tramite acido ed enzimi, processo meccanico mediante contrazioni
muscolari); il pancreas (le cellule esocrine secernono tamponi ed enzimi
digestivi, le cellule endocrine secernono ormoni); il fegato (secrezione
della bile, importante per la digestione dei lipidi, deposito di nutrienti
e numerose altre funzioni vitali); la cistifellea (deposito e
concentrazione della bile); l'intestino tenue (digestione enzimatica e
assorbimento di acqua, substrati organici e ioni); l'intestino crasso
(disidratazione e consolidamento dei materiali non digeriti in
preparazione all'eliminazione).

% ---------------------------------------------------------------------------
\section{La cavità orale}

La \textbf{cavità orale} è situata nello splancnocranio e rappresenta la
prima porzione del canale alimentare. È collocata nella testa, a livello
del massiccio facciale, e comunica con l'esterno tramite la rima buccale,
delimitata dal bordo libero delle labbra. Posteriormente la cavità orale
prosegue nel cosiddetto istmo delle fauci e quindi nella faringe, la quale
si suddivide a sua volta in tre porzioni: la \textbf{rinofaringe},
comunicante anteriormente con le cavità nasali attraverso le coane; la
\textbf{orofaringe}, in continuità con la cavità orale; e la
\textbf{laringofaringe}, che si continua inferiormente con la laringe e
con l'esofago.

La cavità orale comprende denti, gengive, lingua, palato e alcune
formazioni ghiandolari. Denti e gengive si organizzano in due arcate
gengivodentali, una superiore o mascellare e una inferiore o mandibolare.
Il vestibolo della bocca ha forma di ferro di cavallo; le labbra e le
guance presentano una faccia esterna rivestita da cute e una faccia
interna rivestita da tonaca mucosa, separate da muscoli e tessuto
connettivo. Le arcate gengivo-dentali sono costituite dai processi
alveolari della mascella, superiormente, e della mandibola, inferiormente,
e dai denti che in essi si inseriscono.

% ---------------------------------------------------------------------------
\subsection{La lingua}

La lingua è un organo muscolomembranoso che si estende fino alla faringe e
partecipa alla formazione del pavimento della cavità orale propriamente
detta. Le sue funzioni comprendono: il trattamento meccanico del cibo
tramite compressione, abrasione e rimescolamento; l'assistenza alla
preparazione del materiale che verrà deglutito; l'analisi sensoriale
tramite recettori tattili, termici e gustativi; la secrezione di enzimi e
mucine che aiutano la digestione dei grassi.

La lingua si suddivide in radice, corpo e apice. Sul dorso della lingua
un solco mediano la percorre longitudinalmente, mentre il solco terminale,
a forma di V, separa il corpo dalla radice e individua al suo vertice il
foro cieco. Sul corpo della lingua si distribuiscono le papille linguali:
le papille filiformi e le papille fungiformi, disseminate sulla
superficie, e le papille foliate e le papille vallate, disposte in
prossimità del solco terminale. A livello della radice si trova la
tonsilla linguale, costituita da follicoli linguali. Posteriormente, tra
la radice della lingua e l'epiglottide, si individuano la piega
glosso-epiglottica mediana e le pieghe glosso-epiglottiche laterali, che
delimitano le vallecole; ai lati sono presenti gli archi e i muscoli
palato-glossi e palato-faringei, in rapporto con la tonsilla palatina.

% ---------------------------------------------------------------------------
\subsection{I denti}

Durante lo sviluppo si formano due ordini di denti. I primi a comparire
sono i \textbf{denti decidui} (i cosiddetti denti da latte), che spuntano
tra i 6 e i 30 mesi di vita, a cominciare dagli incisivi. Tra i 6 e i 25
anni di età questi vengono sostituiti da 32 denti permanenti, organizzati
secondo una numerazione dentale che distingue, in ciascuna arcata,
un'emiarcata destra e una sinistra: nell'arcata superiore o mascellare i
denti sono numerati da 1 a 16 procedendo da destra a sinistra, mentre
nell'arcata inferiore o mandibolare la numerazione prosegue da 17 a 32
procedendo da sinistra a destra. Per tipologia i denti si distinguono in
incisivi, canini, premolari e molari.

Il dente è costituito da tre parti: la \textbf{corona}, la parte
sporgente all'esterno; la \textbf{radice}, da una a tre, la parte infissa
nell'alveolo; il \textbf{colletto}, la parte intermedia fra le due
precedenti. All'interno del dente si trova la cavità pulpare, occupata
dalla polpa del dente, un tessuto connettivo riccamente vascolarizzato e
innervato, che comunica con l'esterno del dente tramite i canali
radicolari. Ciascuna radice è legata alle pareti del proprio alveolo dal
\textbf{parodonto}, costituito più superficialmente dalla gengiva e, più
profondamente, dal periostio alveolare e da un robusto legamento
connettivale che unisce il periostio al cemento.

I tessuti che compongono il dente sono: lo \textbf{smalto}, che copre la
corona ed è costituito da uno strato calcificato e da matrice organica
ricca di sali di fluoro, prodotta da cellule di origine epiteliale; la
\textbf{dentina}, che forma la massa del dente ed è attraversata dalla
cavità pulpare, le cui cellule (gli odontoblasti) tappezzano la cavità
pulpare inviando lunghi prolungamenti entro i canalicoli che percorrono la
dentina stessa; il \textbf{cemento}, che ricopre la maggior parte della
radice ed è un tessuto osseo particolare le cui cellule specializzate, i
cementociti, sono più numerose a livello dell'apice.

% ---------------------------------------------------------------------------
\subsection{Il palato}

Il palato è suddiviso in due porzioni: il \textbf{palato osseo o duro},
anteriore, composto dai processi palatini delle mascelle e dalle lamine
orizzontali delle ossa palatine, rivestiti da tonaca mucosa; e il
\textbf{palato molle o velo palatino}, posteriore, costituito da una
lamina connettivale rivestita da muscoli (tensore del velo palatino,
elevatore del velo palatino, palatoglosso, palato-faringeo, muscolo
dell'ugola) e da una tonaca mucosa. Centralmente il palato molle si
prolunga verso il basso formando l'\textbf{ugola palatina}, o velo
pendulo. Nella regione del palato decorrono inoltre le ghiandole
palatine, l'arteria e il nervo palatino maggiore, le arterie e i nervi
palatini minori e i rami tonsillari dei nervi palatini minori e del nervo
glossofaringeo, mentre il nervo linguale accompagna il pavimento della
cavità orale; anteriormente le labbra presentano un frenulo labiale
superiore e uno inferiore.

% ---------------------------------------------------------------------------
\section{Le ghiandole salivari}

Le ghiandole salivari maggiori sono tre paia: la \textbf{parotide}, la
\textbf{sottomandibolare} e la \textbf{sottolinguale}. Vi sono inoltre
numerose ghiandole salivari minori, intraparietali, sistemate nella
tonaca propria o nella sottomucosa della cavità orale. Nel loro insieme le
ghiandole salivari producono circa 1--1,5 litri di secreto al giorno.

Le secrezioni salivari sono sotto il controllo del sistema nervoso
autonomo, che le innerva con una doppia innervazione. Qualunque sostanza
presente nella bocca può dare il via a un riflesso salivare tramite
stimolazione dei recettori del nervo trigemino o dei calici gustativi,
innervati dai nervi encefalici VII (faciale), IX (glossofaringeo) e X
(vago).

% ---------------------------------------------------------------------------
\subsection{Parotide}

La parotide è la ghiandola salivare di maggiori dimensioni, con un peso
medio di circa 30 grammi; ha aspetto lobulato, è ricoperta da una robusta
capsula fibrosa e contiene frequentemente tessuto adiposo. Si trova appena
sotto la pelle, anteriormente al lobo dell'orecchio, nella loggia
parotidea del collo, davanti al muscolo sternocleidomastoideo. Produce un
secreto sieroso denso contenente l'amilasi salivare, enzima digestivo che
dà inizio alla digestione chimica dei carboidrati complessi.

% ---------------------------------------------------------------------------
\subsection{Sottomandibolare}

La ghiandola sottomandibolare si trova a metà del corpo della mandibola,
medialmente al suo margine, appena sotto il muscolo miloioideo. Ha
secrezione mista sieromucosa a prevalenza sierosa: gli adenomeri misti
hanno struttura tubulare con lume ampio, ed è ben rappresentata la
componente mioepiteliale. La parte terminale degli adenomeri è chiusa
dalle semilune sierose, la cui secrezione facilita il drenaggio della
componente mucosa.

% ---------------------------------------------------------------------------
\subsection{Sottolinguale}

La ghiandola sottolinguale è la più piccola delle ghiandole salivari
maggiori, con un peso di 3--6 grammi e forma ovoidale. Si trova nel
pavimento della bocca ed è costituita da un gruppetto di piccoli lobuli,
una quindicina, situati al di sotto della mucosa del pavimento della
cavità buccale, in rapporto con i muscoli genioglosso, genioioideo e
ventre anteriore del digastrico, con il muscolo miloioideo, con il
muscolo pterigoideo mediale e con l'osso ioide. Ogni lobulo è provvisto di
un proprio dotto escretore che sbocca direttamente nel solco sublinguale,
in prossimità della caruncola sublinguale, oppure confluisce in un dotto
maggiore comune, il dotto sublinguale maggiore (o dotto di Bartolini); il
dotto della ghiandola sottomandibolare decorre nella stessa regione,
sboccando anch'esso in prossimità della caruncola sublinguale.

% ---------------------------------------------------------------------------
\subsection{La saliva}

La saliva è un biofluido costituito per il 99\% da acqua e per il
restante 1\% da sostanze organiche e inorganiche: elettroliti (sodio,
potassio, magnesio, calcio), muco, sostanze antibatteriche (come le
immunoglobuline IgA), ormoni (tra cui pepsina e testosterone), enzimi
(come l'alfa-amilasi e la lipasi), oltre a cellule, batteri e virus.

% ---------------------------------------------------------------------------
\section{Epiglottide, faringe e deglutizione}

L'\textbf{epiglottide} è una membrana cartilaginea che chiude le vie
aeree nel momento della deglutizione.

I muscoli della faringe coinvolti nella deglutizione sono: i muscoli
costrittori della faringe, che spingono il bolo verso l'esofago; i
muscoli palato-faringeo e stilo-faringeo, che innalzano la laringe; i
muscoli palatali, che innalzano il palato molle e le porzioni adiacenti
della parete faringea.

% ---------------------------------------------------------------------------
\subsection{Le fasi della deglutizione}

La deglutizione si articola in tre fasi successive.

Nella \textbf{fase buccale} il bolo viene spinto contro il palato duro; la
retrazione della lingua spinge il bolo verso la faringe e provoca
l'innalzamento del palato molle. Questa fase è sotto il controllo della
volontà, ma non appena il bolo entra nell'orofaringe si attiva un
riflesso involontario che ne determina il movimento verso lo stomaco.

Nella \textbf{fase faringea}, che inizia quando il bolo prende contatto
con gli archi palatali e con la parete posteriore della faringe,
l'innalzamento della laringe e il ripiegamento dell'epiglottide dirigono
il bolo oltre la glottide, che a questo punto risulta chiusa;
successivamente i muscoli costrittori della faringe lo spingono
nell'esofago.

Nella \textbf{fase esofagea}, che inizia con l'apertura dello sfintere
esofageo superiore, il bolo, appena attraversato lo sfintere, viene spinto
lungo l'esofago da onde peristaltiche; l'avvicinamento del bolo provoca
l'apertura dello sfintere esofageo inferiore, attraverso cui il bolo
passerà nello stomaco.

% ---------------------------------------------------------------------------
\section{L'esofago}

L'esofago è un tubo che collega la faringe allo stomaco; attraverso
movimenti peristaltici fa scendere il bolo allo stomaco. La sua
muscolatura è controllata da riflessi viscerali. È un organo cavo che si
estende da C6 a T10, lungo circa 25 centimetri: attraversa il mediastino
posteriore nella cavità toracica, supera il diaframma ed entra in cavità
addominale, aprendosi nello stomaco a livello di T10.

L'esofago è cilindrico, ma presenta curvature e restringimenti: segue la
cifosi toracica della colonna vertebrale, per cui viene prima spinto
verso destra e poi in avanti a causa dei rapporti con l'arco dell'aorta,
che lo incrocia sulla sinistra, e con l'aorta toracica, che gli decorre
posteriormente. I restringimenti si trovano nel tratto cervicale iniziale,
a livello toracico e in corrispondenza del passaggio attraverso
l'orifizio esofageo del diaframma; sono punti dove il transito del bolo
può essere rallentato. Il drenaggio linfatico dell'esofago e delle
strutture mediastiniche adiacenti coinvolge i linfonodi tracheali, i
linfonodi mediastinici posteriori e parietali posteriori, i linfonodi
tracheo-bronchiali e intercostali, i linfonodi diaframmatici, i linfonodi
gastrici di sinistra (o cardiali dello stomaco), i linfonodi retrocardiaci
e infracardiaci e i linfonodi dell'arteria celiaca, che confluiscono nel
condotto toracico e nei linfonodi giugulari interni.

% ---------------------------------------------------------------------------
\subsection{Nota clinica: il reflusso gastro-esofageo}

Il reflusso, generalmente acido, espone la mucosa dell'esofago all'azione
lesiva dell'acido e degli enzimi contenuti nel succo gastrico. Le cause
sono diverse: il rilasciamento della valvola posta tra esofago e stomaco,
detta \textbf{sfintere esofageo inferiore} (SEI), che normalmente impedisce
al contenuto dello stomaco di tornare nell'esofago, e al cui cattivo
funzionamento possono contribuire cattive abitudini di vita (fumo, caffè,
farmaci) o patologie come l'ernia iatale; il rallentamento dello
svuotamento gastrico, ovvero del passaggio del cibo dallo stomaco
all'intestino, causato ad esempio da pasti abbondanti o da alimenti
particolari come i cibi grassi o il cioccolato.
